% =========================================================
% PARTIE 2 : ALEXANDRE FRANÇOIS
% =========================================================
\section{Efficacité algorithmique : l'approche par collisions}

\begin{frame}{Optimisation : l’approche par collisions}
    
    \begin{block}{Réduction des complexités ($n=64, p=2^{-13}$)}
        \centering
        \renewcommand{\arraystretch}{1.3}
        \small
        \begin{tabular}{lccc}
            \toprule
            \textbf{Algorithme} & \textbf{Formule} & \textbf{Opérations} & \textbf{Faisabilité} \\
            \midrule
            Naïf Exhaustif & $2^{2n}$ & $2^{128}$ & Impraticable \\
            Naïf Optimisé  & $2^n \cdot p^{-1}$ & $2^{77}$ & Supercalculateurs \\
            \midrule
            \textbf{Par Collisions} & $\mathbf{2^{n/2} \cdot p^{-1}}$ & $\mathbf{2^{45}}$ & \textbf{Station de travail} \\
            \bottomrule
        \end{tabular}
    \end{block}

    \vspace{0.4cm}

    \begin{itemize}
        \item \textbf{Principe de collision :} Trouver deux messages distincts $x$ et $x'$ produisant la même signature de sortie (sur une fonction dérivée $g_\gamma$).
        \item \textbf{Paradoxe des anniversaires :} Réduction conjointe de la complexité en temps et en espace.
    \end{itemize}
    
\end{frame}

% --- DIAPO 2 : CONCEPT DE SIGNATURE ---
\begin{frame}{Le concept de signature : la fonction $g_\gamma$}
    \begin{block}{Définition du substitut}
        On choisit une différence arbitraire $\gamma \neq 0$, indépendante de la différence $\alpha$ recherchée.
    \end{block}
    
    \begin{block}{Fonction dérivée $g_\gamma$}
        Pour chaque message $x$, on calcule sa signature de sortie :
        \[ g_{\gamma}(x) = F(x) \oplus F(x \oplus \gamma) \]
    \end{block}

    \begin{itemize}
        \item L'algorithme ne fixe plus $\alpha$ a priori.
        \item On recherche une \textbf{collision globale} : $g_\gamma(x) = g_\gamma(x')$.
        \item La direction $\alpha$ est révélée par la collision : $\alpha = x \oplus x'$.
    \end{itemize}
\end{frame}

% --- DIAPO 3 : PREUVE ALGÉBRIQUE ---
\begin{frame}{Preuve algébrique de la collision}
    \begin{block}{Double différenciation}
        En calculant le XOR des deux signatures, on réorganise les termes :
        \begin{align*}
            g_\gamma(x) \oplus g_\gamma(x \oplus \alpha) &= \big( F(x) \oplus F(x \oplus \gamma) \big) \oplus \big( F(x \oplus \alpha) \oplus F(x \oplus \alpha \oplus \gamma) \big) \\
            &= \big( F(x) \oplus F(x \oplus \alpha) \big) \oplus \big( F(x \oplus \gamma) \oplus F(x \oplus \alpha \oplus \gamma) \big)
        \end{align*}
    \end{block}
    \begin{block}{Annulation de la différence $\beta$}
        Si les deux couples satisfont la différentielle $\alpha \to \beta$ (quadruplet correct) :
        \begin{equation*}
            g_\gamma(x) \oplus g_\gamma(x \oplus \alpha) = \beta \oplus \beta = 0
        \end{equation*}
    \end{block}
    \begin{itemize}
        \item L'égalité $g_\gamma(x) = g_\gamma(x \oplus \alpha)$ est donc garantie.
    \end{itemize}
\end{frame}

% --- DIAPO 4 : DÉTECTION (MÉCANIQUE) ---
\begin{frame}{Phase de détection : la mécanique}
    
    \begin{block}{Changement de paradigme}
        \begin{itemize}
            \item \textbf{Approche naïve :} Itération sur les $2^n$ valeurs de $\alpha$.
            \item \textbf{Par collisions :} Traitement global sur un échantillon de $M$ messages.
        \end{itemize}
    \end{block}

    \vspace{0.4cm}
    \textbf{Algorithme en 3 étapes clés :}
    \vspace{0.2cm}

    \begin{enumerate}
        \item \textbf{Générer :} Calculer les signatures $g_\gamma(x)$ pour $M$ messages.
        \item \textbf{Hacher :} Détecter les doublons ($g_\gamma(x) = g_\gamma(x')$).
        \item \textbf{Déduire :} Isoler les candidats $\alpha = x \oplus x'$.
    \end{enumerate}

\end{frame}

% --- DIAPO 5 : DIMENSIONNEMENT ET VÉRIFICATION ---
\begin{frame}{Dimensionnement par le paradoxe des anniversaires}

    \begin{block}{Espérance des collisions utiles}
        Pour un échantillon de taille $M$ :
        \[ E[\text{Collisions}] = \frac{M^2}{2} \cdot 2^{-n} \cdot p^2 \]
    \end{block}

    \begin{itemize}
        \item $\mathbf{2^{-n}}$ : Probabilité d'obtenir la différence ciblée $\alpha$.
        \item $\mathbf{p^2}$ : Probabilité d'obtenir un quadruplet correct.
    \end{itemize}

    \vspace{0.2cm}

    \begin{block}{Taille optimale}
        Pour garantir la détection ($E \ge 1$), l'échantillon requis s'effondre à :
        \[ \mathbf{M \approx 2^{n/2} \cdot p^{-1}} \]
    \end{block}

    \vspace{0.2cm}
    \textit{\small $\rightarrow$ Un test de vérification classique élimine ensuite les éventuels faux positifs.}

\end{frame}